\documentclass[11pt,a4paper]{article}
\usepackage[english]{babel}
\usepackage[utf8]{inputenc}
\usepackage[T1]{fontenc}
\usepackage{amsmath,amssymb,bm}
\usepackage{geometry}
\usepackage{xcolor}
\usepackage{booktabs}
\usepackage{hyperref}
\geometry{margin=2.5cm}
\hypersetup{colorlinks=true,linkcolor=blue,urlcolor=blue}

\newcommand{\IS}{\delta}
\newcommand{\QS}{\Delta E_Q}
\newcommand{\BHF}{B_{\mathrm{HF}}}
\newcommand{\dd}{\mathrm{d}}

\title{Isomer shift distributions $P(\IS)$ and distributions $P(\IS,\QS)$ in M\"{o}ssbauer spectra}
\author{Fitbauer -- mathematical implementation note}
\date{\today}

\begin{document}
\maketitle

\section{Physical motivation}

The isomer shift $\IS$ measures the displacement of the centre of the
M\"{o}ssbauer pattern with respect to a reference, normally $\alpha$-Fe. It is
related to the electron density at the nucleus and therefore to the oxidation
state, coordination, and covalency of the Fe atom. In ordered materials a single
$\IS$ per phase is usually fitted. However, in glasses, silicates, amorphous
materials, ferritins, nanoparticles, or mixtures of local environments, a
continuous distribution of shifts may exist.

Fitbauer implements two extensions:
\begin{itemize}
  \item a one-dimensional distribution $P(\IS)$;
  \item two-dimensional distributions combining $\IS$ with another quantity,
  especially $P(\IS,\QS)$ and $P(\BHF,\IS)$.
\end{itemize}
The case $P(\IS,\QS)$ is particularly useful for broad paramagnetic doublets,
where the centre of the doublet and its splitting can vary simultaneously.

\section{One-dimensional distribution $P(\IS)$}

For a grid of shifts $\IS_i$, the 1D model is
\begin{equation}
  T(v_k)=b_0+b_1v_k-\sum_i P_i\,S(v_k;\IS_i,\QS,\BHF,\Gamma),
\end{equation}
where $S$ is the unit-depth absorption pattern. If $\BHF=0$, the pattern reduces
in practice to a doublet or singlet depending on the value of $\QS$.
The weights satisfy
\begin{equation}
  P_i\ge0.
\end{equation}
The regularised problem is written as
\begin{equation}
  \Phi(\bm p)=\|W(\bm y-X\bm p)\|^2+\alpha\|D^2\bm p\|^2,
\end{equation}
where $D^2$ is the second-difference operator. The normalised probability is
computed as
\begin{equation}
  p(\IS_i)=\frac{P_i}{\int P(\IS)\,\dd\IS}.
\end{equation}

\section{Two-dimensional distribution $P(\IS,\QS)$}

For a grid $\IS_i$, $Q_j$, the model is
\begin{equation}
  T(v_k)=b_0+b_1v_k-\sum_{ij}P_{ij}\,
  S(v_k;\IS_i,Q_j,\BHF,\Gamma).
\end{equation}
To describe a paramagnetic population, $\BHF=0$ is fixed and each point
$(\IS_i,Q_j)$ generates a doublet centred at $\IS_i$ with splitting $Q_j$:
\begin{align}
  v_{1,ij}&=\IS_i-\frac{Q_j}{2},\\
  v_{2,ij}&=\IS_i+\frac{Q_j}{2}.
\end{align}

The 2D regularisation penalises curvature in both directions:
\begin{equation}
  \Phi(\bm P)=\|W(\bm y-X\bm z)\|^2
  +\alpha_{\IS}\|D_{\IS}^2P\|^2
  +\alpha_Q\|D_Q^2P\|^2,
  \qquad P_{ij}\ge0.
\end{equation}
In augmented form:
\begin{equation}
  \begin{bmatrix}
    WX\\
    \sqrt{\alpha_{\IS}}\,D_{\IS}^2\\
    \sqrt{\alpha_Q}\,D_Q^2
  \end{bmatrix}\bm z\simeq
  \begin{bmatrix}
    W\bm y\\
    \bm 0\\
    \bm 0
  \end{bmatrix}.
\end{equation}

\section{Interpretation of $P(\IS,\QS)$}

The marginals are
\begin{align}
  P_{\IS}(\IS_i)&=\int P(\IS_i,Q)\,\dd Q,\\
  P_Q(Q_j)&=\int P(\IS,Q_j)\,\dd\IS.
\end{align}
A tilted cloud in the map may indicate that sites with higher $\IS$ also have
higher or lower $\QS$. This may suggest families of local environments, for
example variations in valence, coordination, or polyhedral distortion.

Fitbauer computes descriptive diagnostics of the map:
\begin{align}
  \langle\IS\rangle,\quad \sigma_{\IS},\quad
  \langle Q\rangle,\quad \sigma_Q,\quad
  \rho_{\IS Q}.
\end{align}
The apparent correlation is
\begin{equation}
  \rho_{\IS Q}=\frac{\langle(\IS-\langle\IS\rangle)(Q-\langle Q\rangle)\rangle}
  {\sigma_{\IS}\sigma_Q}.
\end{equation}
These values describe the fitted distribution but do not by themselves prove
that the distribution is unique or physical.

\section{Distribution $P(\BHF,\IS)$}

$P(\BHF,\IS)$ is also implemented:
\begin{equation}
  T(v_k)=b_0+b_1v_k-\sum_{ij}P_{ij}\,
  S(v_k;B_i,\QS,\IS_j,\Gamma).
\end{equation}
It can be useful in magnetic phases where the hyperfine field and the isomer
shift vary jointly, for example due to size distributions, partial oxidation,
or surface environments. However, it is especially sensitive to velocity
calibration errors: a global shift of the velocity scale can mimic an $\IS$
distribution.

\section{Identifiability risks}

\textbf{Important warning:} distributions involving $\IS$ are very sensitive to
calibration, folding, and background. In a doublet,
\begin{equation}
  v_1=\IS-\frac{\QS}{2},\qquad v_2=\IS+\frac{\QS}{2},
\end{equation}
so changes in $\IS$ shift the centre and changes in $\QS$ modify the splitting.
With broad lines or noise, many combinations $(\IS,\QS)$ can produce similar
spectra. Moreover, $\IS$ correlates with:
\begin{itemize}
  \item the folding point;
  \item the calibration $V_{\max}$;
  \item the isomeric reference;
  \item the linewidth $\Gamma$;
  \item the slope/background;
  \item omitted discrete components.
\end{itemize}
Therefore, a visually smooth $P(\IS)$ or $P(\IS,\QS)$ with a good residual may
lack physical meaning if it is compensating an instrumental error or an
incomplete model.

\section{Best practices}

Before interpreting a distribution involving $\IS$ it is recommended to:
\begin{enumerate}
  \item carefully fix or verify the velocity calibration;
  \item use a known isomeric reference;
  \item check the folding point with the residual diagnostic;
  \item compare with discrete models and with 1D $P(\QS)$;
  \item vary $\alpha_{\IS}$ and $\alpha_Q$ over several orders of magnitude;
  \item repeat with a different number of bins;
  \item avoid publishing fine details that disappear when the regularisation
  is changed;
  \item use external information: chemistry, XRD, TEM, magnetometry, or
  literature.
\end{enumerate}
For publication, it is advisable to show that the $P(\IS,\QS)$ map is stable
and that a simpler model does not explain the data satisfactorily.

\section{Status in Fitbauer}

The GUI is accessed via the modes:
\begin{itemize}
  \item \textbf{P(IS)} for the 1D distribution;
  \item \textbf{P(IS,$\Delta E_Q$) 2D};
  \item \textbf{P(BHF,IS) 2D}.
\end{itemize}
The 2D modes share the regularised engine used for $P(\BHF,\QS)$. The GUI
shows a heatmap, marginals, a regularisation L-surface, TSV export, an
interactive Plotly heatmap, and a PDF report with the 2D map.

\end{document}
